\documentclass[11pt]{amsart}
\usepackage{amsmath,amssymb,amsthm}

\newtheorem{theorem}{Theorem}[section]
\newtheorem{lemma}[theorem]{Lemma}
\newtheorem{proposition}[theorem]{Proposition}
\newtheorem{corollary}[theorem]{Corollary}
\newtheorem{conjecture}[theorem]{Conjecture}
\theoremstyle{definition}
\newtheorem{definition}[theorem]{Definition}
\theoremstyle{remark}
\newtheorem*{remark}{Remark}

\title{Divisibility and coprimality in large subsets of $\{1,\dots,2n\}$}
\author{T. W. Mockup}
\address{Laboratory of Fictitious Combinatorics (fictional)}
\email{t.w.mockup@example.org}
\date{}

\begin{document}

\begin{abstract}
Two classical pigeonhole facts state that any $n+1$ numbers chosen from
$\{1,\dots,2n\}$ contain two coprime numbers, and contain two numbers one of
which divides the other. We give short proofs of both, show that they are
sharp, and compare the second with the Erd\H{o}s--Szekeres theorem on
monotone subsequences.
\end{abstract}

\maketitle

\section{Introduction}

The divisibility statement (Theorem~\ref{thm:main}) is a well-known problem
of Erd\H{o}s; it is discussed together with several other pigeonhole
arguments in~\cite{AZ}. We give a self-contained treatment in which the role
of the odd part of an integer is made explicit, and we record the matching
extremal examples (Propositions~\ref{prop:sharpdiv} and~\ref{prop:sharpcop}).

\section{Definitions and the two pigeonhole theorems}

\begin{definition}\label{def:basic}
For an integer $m\ge 0$ we write $[m]=\{1,2,\dots,m\}$. A set $A$ of positive
integers is a \emph{divisibility antichain} if there are no $a,b\in A$ with
$a\mid b$, and a \emph{divisibility chain} if for all $a,b\in A$ we have
$a\mid b$ or $b\mid a$.
\end{definition}

\begin{lemma}\label{lem:consec}
For every integer $k\ge 1$ we have $\gcd(k,k+1)=1$.
\end{lemma}

\begin{proof}
Any common divisor of $k$ and $k+1$ divides $(k+1)-k=1$.
\end{proof}

\begin{proposition}\label{prop:coprime}
Let $n\ge 1$ and let $A\subseteq[2n]$ with $|A|\ge n$. Then $A$ contains two
distinct elements that are coprime.
\end{proposition}

\begin{proof}
Partition $[2n]$ into the $n$ pairs $\{1,2\},\{3,4\},\dots,\{2n-1,2n\}$. Since
the $n+1$ (or more) elements of $A$ are distributed among $n$ pairs, two of
them lie in the same pair, i.e.\ they are consecutive integers $k,k+1$, and
these are coprime by Lemma~\ref{lem:consec}.
\end{proof}

\begin{lemma}\label{lem:oddpart}
Every integer $m\ge 1$ can be written uniquely as $m=2^{k}q$ with $k\ge 0$ and
$q$ odd. We call $q$ the \emph{odd part} of $m$.
\end{lemma}

\begin{proof}
Take $k$ to be the exponent of $2$ in the prime factorisation of $m$, and
$q=m/2^k$. Uniqueness follows from uniqueness of prime factorisation.
\end{proof}

\begin{theorem}[Main theorem]\label{thm:main}
Let $n\ge 1$ and let $A\subseteq[2n]$ with $|A|\ge n+1$. Then there are
distinct $a,b\in A$ with $a\mid b$.
\end{theorem}

\begin{proof}
The odd part of an element of $[2n]$ is an odd number in $[2n]$, so it lies
in $\{1,3,\dots,2n-1\}$, a set of $n$ elements. By the pigeonhole principle
two distinct elements of $A$ have the same odd part $q$, say $a=2^{i}q$ and
$b=2^{j}q$ with $i\ne j$; we may assume $i<j$. Then $b\mid a$, since
$a=2^{i-j}b$.
\end{proof}

\section{Sharpness}

\begin{proposition}\label{prop:sharpdiv}
For every $n\ge 1$ the set $\{n+1,n+2,\dots,2n\}\subseteq[2n]$ is a
divisibility antichain with $n$ elements. Hence the bound $n+1$ in
Theorem~\ref{thm:main} cannot be lowered.
\end{proposition}

\begin{proof}
If $a\ne b$ are in the set and $a\mid b$, then $b\ge 2a\ge 2(n+1)>2n$, which is
impossible.
\end{proof}

\begin{proposition}\label{prop:sharpcop}
For every $n\ge 1$ the set $\{2,4,\dots,2n\}\subseteq[2n]$ has $n$ elements,
no two of which are coprime.
\end{proposition}

\begin{proof}
Any two of its elements have the common divisor $2$.
\end{proof}

\begin{lemma}\label{lem:nocop}
Let $n\ge 1$ and let $A\subseteq[2n]$ contain no two distinct coprime
elements. Then $|A|\le n$, and every element of $A$ is even.
\end{lemma}

\begin{proof}
The bound $|A|\le n$ is the contrapositive of Proposition~\ref{prop:coprime}.
If $a\in A$ were odd, then $a$ would be coprime to both $a-1$ and $a+1$ by
Lemma~\ref{lem:consec}, and by maximality one of these lies in $A$, a
contradiction.
\end{proof}

\begin{remark}
Theorem~\ref{thm:main} says that $[2n]$ contains no divisibility antichain of
size $n+1$. It is instructive to compare this with the theorem of
Erd\H{o}s and Szekeres~\cite{ES}: \emph{every sequence of $n^2$ distinct real
numbers contains a monotone subsequence of length $n+1$.} Both are
consequences of the pigeonhole principle applied to a suitable labelling.
\end{remark}

\begin{thebibliography}{9}
\bibitem{AZ} M.~Aigner and G.~M.~Ziegler, \emph{Proofs from THE BOOK},
6th ed., Springer, 2018.
\bibitem{ES} P.~Erd\H{o}s and G.~Szekeres, \emph{A combinatorial problem in
geometry}, Compositio Math. \textbf{2} (1935), 463--470.
\end{thebibliography}

\end{document}
